\chapter{Creio em Jesus Cristo, seu único filho, Nosso Senhor}
\chaptermark{Creio em Jesus Cristo}

O segundo artigo da Profissão de Fé, ``Creio em Jesus Cristo, seu único Filho, nosso Senhor'', coloca no centro da fé cristã a pessoa de Jesus de Nazaré. Este artigo aprofunda o mistério da identidade de Jesus, revelando o significado do seu nome, dos títulos que Lhe são atribuídos e da sua relação única com Deus Pai. Baseado nos parágrafos 430 a 455 do Catecismo da Igreja Católica, o texto que se segue apresenta de forma acessível e ampliada os principais ensinamentos sobre o nome ``Jesus'', o título ``Cristo'', a filiação divina e a soberania do Senhor. Estas verdades não são meras informações teológicas, mas o coração da experiência cristã, que nos convida a entrar numa relação pessoal de fé, amor e adesão a Jesus, o Salvador do mundo.

\section{O Nome ``Jesus'' (430--435)}

O nome ``Jesus'' significa, em hebraico, ``Deus salva''. Este nome não foi escolhido pelos homens, mas dado pelo próprio Deus por meio do anjo Gabriel na anunciação. Desde o início da sua vida terrena, Jesus recebe um nome que exprime a sua identidade mais profunda e a sua missão: salvar o seu povo dos seus pecados.

Deus, ao longo da história da salvação, manifestou-se como aquele que liberta o seu povo. Primeiro libertou Israel da escravidão do Egipto. Depois, progressivamente, revelou que o verdadeiro perigo é o pecado, que ofende o próprio Deus. Só Deus pode perdoar o pecado. Por isso, o nome de Jesus revela que n’Ele o próprio Deus se fez presente para salvar a humanidade de forma definitiva.

O nome de Jesus é, portanto, o nome do Salvador. Nele está presente o próprio nome divino. Todos os que invocam este nome com fé podem receber a salvação. A Igreja, desde os primeiros tempos, reconheceu que ``não há outro nome debaixo do céu pelo qual possamos ser salvos''. Invocar o nome de Jesus é entrar em contato com o poder salvador de Deus feito homem.

Este nome tão doce e poderoso deve estar sempre nos lábios e no coração dos fiéis. Na oração, na tentação, na alegria e no sofrimento, o nome de Jesus é fonte de força e de paz. Ele resume toda a esperança cristã: Deus não nos abandonou no pecado, mas veio salvar-nos pessoalmente.

O nome de Jesus convida-nos a uma confiança total. Não se trata apenas de um nome histórico, mas de uma presença viva que continua a agir na Igreja através dos sacramentos e da oração. Quem crê em Jesus crê no Deus que salva.

\section{O Título ``Cristo'' (436--440)}

O título ``Cristo'' significa ``Ungido''. No Antigo Testamento, eram ungidos os reis, os sacerdotes e os profetas para uma missão especial da parte de Deus. Jesus reúne em Si mesmo estas três dimensões de forma perfeita e definitiva. Ele é o Ungido por excelência, enviado pelo Pai para cumprir a missão da salvação.

Desde o início do seu ministério, Jesus é reconhecido como o Messias esperado por Israel. Os apóstolos, especialmente Pedro, confessam publicamente: ``Tu és o Cristo, o Filho do Deus vivo''. Esta confissão é o fundamento da fé da Igreja. Jesus não é um profeta entre outros, mas o próprio cumprimento de todas as promessas de Deus.

O título de Cristo revela que Jesus é o centro da história da salvação. Tudo o que aconteceu antes d’Ele apontava para Ele. Depois d’Ele, a história da humanidade é iluminada pela sua luz. A unção de Jesus é a unção do Espírito Santo, que desce sobre Ele no Batismo e no Pentecostes se derrama sobre a Igreja.

Ser cristão significa, literalmente, ser ``outro Cristo''. Os que recebem o Batismo são ungidos com o mesmo Espírito e chamados a participar na missão profética, sacerdotal e real de Jesus. O título de Cristo transforma a nossa identidade: já não pertencemos a nós mesmos, mas a Aquele que nos ungiu para uma vida nova.

A fé em Jesus Cristo exige que O reconheçamos como o Ungido de Deus e que nos deixemos transformar pela sua unção. Esta fé não é abstrata, mas concreta: implica seguir os seus passos e anunciar o Evangelho com a vida.

\section{Seu Único Filho (441--445)}

Jesus é confessado como o ``único Filho'' de Deus. Esta filiação não é adotiva como a nossa, mas eterna e única. O Pai e o Filho são uma só substância divina. Jesus revela o Pai como ninguém jamais pôde fazê-lo: ``Quem me vê a mim vê o Pai''.

A confissão da filiação divina de Jesus é o ponto alto da revelação. Nos evangelhos, o próprio Deus declara no Batismo e na Transfiguração: ``Este é o meu Filho amado''. Esta verdade ultrapassa todas as expectativas humanas e só pode ser acolhida pela fé. A Igreja, guiada pelo Espírito Santo, professou sempre esta filiação divina contra todas as heresias.

O fato de Jesus ser o Filho único manifesta o amor infinito de Deus pela humanidade. O Pai não enviou um mensageiro, mas o seu próprio Filho. Na Encarnação, o Filho de Deus assume a nossa humanidade para nos tornar filhos adotivos do Pai. Esta é a maravilha da nossa fé.

A filiação divina de Jesus ilumina todo o mistério cristão. A Cruz, a Ressurreição e a Ascensão só podem ser compreendidas à luz desta relação única entre o Pai e o Filho. A Eucaristia, por sua vez, é o dom que o Filho faz de Si mesmo ao Pai e a nós.

Reconhecer Jesus como o Filho único de Deus transforma a nossa oração. Já não oramos como estranhos, mas como filhos que se dirigem ao Pai com a confiança que o próprio Jesus nos ensinou.

\section{Nosso Senhor (446--451)}

O título ``Senhor'' aplicado a Jesus é uma das confissões mais antigas da Igreja. Na linguagem bíblica, ``Senhor'' era o nome reservado ao próprio Deus. Ao chamar Jesus de ``Senhor'', a Igreja primitiva reconhecia a sua divindade. ``Jesus é o Senhor'' era o grito de fé dos primeiros cristãos.

Este título manifesta a soberania de Jesus sobre o universo e sobre a história. Ele não é um senhor entre outros, mas o Senhor único. Diante d’Ele, todo o joelho se dobra no céu, na terra e debaixo da terra. A sua senhoria não é opressiva, mas libertadora, pois Ele nos liberta do pecado e da morte.

Professar Jesus como ``nosso'' Senhor implica uma entrega pessoal. Não basta reconhecê-Lo como Senhor do universo; é necessário aceitá-Lo como Senhor da nossa vida concreta. Esta confissão exige obediência, confiança e amor total.

A soberania de Cristo estende-se também à Igreja. Ele é a Cabeça do Corpo místico. Os sacramentos, especialmente a Eucaristia, são os meios pelos quais o Senhor ressuscitado continua a estar presente e a agir no meio do seu povo.

A confissão de Jesus como Senhor dá sentido a toda a existência cristã. Nas alegrias e nas tribulações, o cristão sabe que pertence Àquele que venceu o mundo. Esta certeza enche o coração de paz e de coragem missionária.

\section{A Fé no Senhor Jesus (452--455)}

A fé em Jesus Cristo, Filho único e Senhor, é o centro da vida cristã. Esta fé não é uma opinião entre muitas, mas a resposta ao amor de Deus que se revelou plenamente n’Ele. A Igreja, ao longo dos séculos, guardou e transmitiu esta fé, defendendo-a contra os erros que diminuem a divindade ou a humanidade de Jesus.

Esta profissão de fé tem consequências práticas para a vida diária. Quem crê em Jesus como Senhor procura conformar a sua vida aos seus ensinamentos e exemplos. A caridade, o perdão, a justiça e a pureza tornam-se expressões concretas desta fé.

A fé em Cristo ilumina também o sentido do sofrimento e da morte. Porque Jesus é o Senhor que morreu e ressuscitou, podemos enfrentar o mistério do mal com esperança. Nada pode separar-nos do amor de Cristo.

Finalmente, esta fé orienta o olhar para o futuro. Jesus, que é o Senhor da história, virá novamente em glória. A sua vinda final será a manifestação plena de tudo o que cremos n’Ele agora pela fé.

A adesão fiel a Jesus Cristo, seu único Filho e nosso Senhor, é o caminho para a verdadeira felicidade e para a vida eterna. Que o Espírito Santo fortaleça em nós esta fé e nos torne testemunhas alegres do Evangelho.

\section{Prefigurações no Antigo Testamento}

O Antigo Testamento contém inúmeras prefigurações de Jesus Cristo. O nome de Jesus evoca Josué, que conduziu o povo à Terra Prometida, tal como Jesus nos conduz à verdadeira Terra Prometida que é o Reino dos Céus. Os ungidos do Antigo Testamento (reis, sacerdotes e profetas) apontam para Aquele que seria o Ungido por excelência.

A figura do Servo Sofredor descrito por Isaías prefigura o mistério da Paixão de Cristo. O sacrifício de Isaac por Abraão aponta para o sacrifício do Filho único que o Pai não poupou por nós. Estas figuras mostram a unidade do plano salvífico de Deus.

A espera do Messias percorre todo o Antigo Testamento como um fio condutor. Os profetas anunciaram o seu nascimento da Virgem, o seu reinado de justiça e a sua missão salvadora. Jesus cumpre todas estas promessas de modo surpreendente e definitivo.

\section{A Promessa de Salvação em Jesus Cristo}

Toda a história da salvação converge para Jesus Cristo. O nome que significa ``Deus salva'' revela que a promessa feita no Paraíso se cumpre n’Ele. Através da sua Encarnação, Paixão, Morte e Ressurreição, Jesus realiza a salvação que Deus preparou desde a eternidade.

A fé em Jesus como o único Salvador exclui qualquer outro caminho de salvação. Ele é o caminho, a verdade e a vida. Quem n’Ele crê passa da morte para a vida. Esta verdade deve encher-nos de gratidão e de zelo missionário.

A Igreja continua a anunciar ao mundo que Jesus Cristo é o Senhor e que só n’Ele há salvação plena. Esta mensagem, antiga e sempre nova, é a maior esperança da humanidade.
